% Options for packages loaded elsewhere
\PassOptionsToPackage{unicode}{hyperref}
\PassOptionsToPackage{hyphens}{url}
\documentclass[
  ignorenonframetext,
]{beamer}
\newif\ifbibliography
\usepackage{pgfpages}
\setbeamertemplate{caption}[numbered]
\setbeamertemplate{caption label separator}{: }
\setbeamercolor{caption name}{fg=normal text.fg}
\beamertemplatenavigationsymbolsempty
% remove section numbering
\setbeamertemplate{part page}{
  \centering
  \begin{beamercolorbox}[sep=16pt,center]{part title}
    \usebeamerfont{part title}\insertpart\par
  \end{beamercolorbox}
}
\setbeamertemplate{section page}{
  \centering
  \begin{beamercolorbox}[sep=12pt,center]{section title}
    \usebeamerfont{section title}\insertsection\par
  \end{beamercolorbox}
}
\setbeamertemplate{subsection page}{
  \centering
  \begin{beamercolorbox}[sep=8pt,center]{subsection title}
    \usebeamerfont{subsection title}\insertsubsection\par
  \end{beamercolorbox}
}
% Prevent slide breaks in the middle of a paragraph
\widowpenalties 1 10000
\raggedbottom
\AtBeginPart{
  \frame{\partpage}
}
\AtBeginSection{
  \ifbibliography
  \else
    \frame{\sectionpage}
  \fi
}
\AtBeginSubsection{
  \frame{\subsectionpage}
}
\usepackage{iftex}
\ifPDFTeX
  \usepackage[T1]{fontenc}
  \usepackage[utf8]{inputenc}
  \usepackage{textcomp} % provide euro and other symbols
\else % if luatex or xetex
  \usepackage{unicode-math} % this also loads fontspec
  \defaultfontfeatures{Scale=MatchLowercase}
  \defaultfontfeatures[\rmfamily]{Ligatures=TeX,Scale=1}
\fi
\usepackage{lmodern}
\ifPDFTeX\else
  % xetex/luatex font selection
\fi
% Use upquote if available, for straight quotes in verbatim environments
\IfFileExists{upquote.sty}{\usepackage{upquote}}{}
\IfFileExists{microtype.sty}{% use microtype if available
  \usepackage[]{microtype}
  \UseMicrotypeSet[protrusion]{basicmath} % disable protrusion for tt fonts
}{}
\makeatletter
\@ifundefined{KOMAClassName}{% if non-KOMA class
  \IfFileExists{parskip.sty}{%
    \usepackage{parskip}
  }{% else
    \setlength{\parindent}{0pt}
    \setlength{\parskip}{6pt plus 2pt minus 1pt}}
}{% if KOMA class
  \KOMAoptions{parskip=half}}
\makeatother
\setlength{\emergencystretch}{3em} % prevent overfull lines
\providecommand{\tightlist}{%
  \setlength{\itemsep}{0pt}\setlength{\parskip}{0pt}}
\usepackage{bookmark}
\IfFileExists{xurl.sty}{\usepackage{xurl}}{} % add URL line breaks if available
\urlstyle{same}
\hypersetup{
  hidelinks,
  pdfcreator={LaTeX via pandoc}}

\author{\texorpdfstring{}{}}
\date{}

\begin{document}

\begin{frame}
\newpage
\end{frame}

\section{Conclusion}\label{sec:conclusion}

\begin{frame}{Summary of Contributions}
\protect\phantomsection\label{summary-of-contributions}
This paper provided comprehensive simulation-based empirical validation
of the Cognitive Integrity Framework (CIF) introduced in Part 1 of this
series. Our contributions span implementation, evaluation, and analysis:

\textbf{Implementation}: We implemented the complete CIF defense
suite---cognitive firewalls, belief sandboxes, trust calculus with
bounded delegation, tripwire detection, behavioral invariants, and
Byzantine-tolerant consensus---as production-ready Python modules with
1,557 passing tests at 100\% pass rate, demonstrating that the formal
mechanisms translate into deployable, independently testable
code.\footnote{Source code available at \url{<https://github.com/docxology/cognitive_integrity}> (DOI: 10.5281/zenodo.18364128)}

\textbf{Attack Corpus}: We assembled 950 cognitive attacks across four
categories (prompt injection, trust exploitation, belief manipulation,
coordination attacks), enabling reproducible security evaluation of
multiagent systems.

\textbf{Cross-Architecture Evaluation}: We evaluated CIF's detection
architecture across six production multiagent topologies (Claude Code,
AutoGPT, CrewAI, LangGraph, MetaGPT, Camel) using parametric
architecture-aware simulation calibrated to published benchmarks. The
simulation models each architecture's topology and attack-surface
exposure to produce detection rates that characterize CIF's design-level
protection properties.

\textbf{Statistical Rigor}: We provided significance testing
(\(p < 0.0001\) for primary hypotheses), effect sizes (Cohen's
\(d > 1.0\) for all major comparisons), confidence intervals, and
ablation studies establishing the robustness of our findings under the
simulation model.
\end{frame}

\begin{frame}{Key Findings}
\protect\phantomsection\label{key-findings}
The simulation-based evaluation yields four principal findings, each
reflecting CIF's design-level detection properties under calibrated
conditions:

\begin{enumerate}
\item
  \textbf{Layered defense is essential}: No single mechanism achieves
  acceptable protection in simulation; composition yields multiplicative
  improvement consistent with theoretical predictions from the defense
  composition algebra.
\item
  \textbf{Trust calculus prevents amplification}: The \(\delta^d\) decay
  bound successfully prevented trust laundering across all tested
  architectures---a structural guarantee that holds independent of
  attacker sophistication and is verified both formally (Part 1) and
  through unit-tested implementation.
\item
  \textbf{Architecture matters}: Peer-to-peer architectures show
  greatest improvement from CIF in simulation, consistent with Part 1's
  prediction that equal-trust topologies are most vulnerable to lateral
  movement attacks.
\item
  \textbf{Performance overhead is manageable}: 20-25\% estimated latency
  overhead for full CIF deployment was observed in simulation, with
  overhead dominated by the cognitive firewall and Byzantine consensus
  components.
\end{enumerate}
\end{frame}

\begin{frame}{Observed Deployment Properties}
\protect\phantomsection\label{observed-deployment-properties}
The evaluation data establishes four empirical properties relevant to
deployment:

\begin{enumerate}
\item
  \textbf{Layered defense is necessary for high efficacy}: No single
  mechanism exceeded 85\% detection. The Minimal-C configuration
  (Firewall + Tripwires + Drift) achieved 90\% at 12\% overhead; full
  CIF reached 94\% at 20--25\% overhead (\cref{tab:minimal-configs}).
\item
  \textbf{Defense efficacy is architecture-dependent}: Tripwire-only
  deployments achieved 82\% detection in hierarchical topologies but
  only 61\% in peer-to-peer systems. Trust calculus with
  \(\delta \leq 0.8\) was the dominant factor in peer-to-peer defense
  (\cref{tab:architecture-insights}).
\item
  \textbf{Detection degrades against novel attacks}: Cross-validation
  with held-out attack types showed 4--10\% detection rate gaps, with
  coordination attacks exhibiting the largest generalization gap
  (\(-10\%\)) (\cref{sec:robustness}).
\item
  \textbf{Byzantine tolerance requires \(n \geq 3f + 1\)}: The minimum
  viable configuration for tolerating a single compromised agent
  (\(f = 1\)) is \(n \geq 4\) agents.
\end{enumerate}

Detailed deployment guidance, including configuration checklists and
operational procedures derived from these findings, is provided in Part
3 of this series.
\end{frame}

\begin{frame}{Alignment with Emerging Standards}
\protect\phantomsection\label{alignment-with-emerging-standards}
CIF's design anticipates and directly addresses the security risks
codified by two major 2025--2026 standardization efforts.

The \textbf{OWASP Top 10 for Agentic Applications} (2026) identifies 10
agentic-specific risks (ASI01--ASI10) \cite{owasp2025agentic}. CIF's
defense mechanisms map systematically to these risks: the Cognitive
Firewall counters Agent Goal Hijack (ASI01) by detecting and filtering
prompt injections before they alter agent objectives; the Belief Sandbox
addresses Tool Misuse and Exploitation (ASI02) by isolating unverified
tool outputs before they propagate into the agent's belief state; Trust
Calculus with \(\delta^d\) decay prevents Identity and Privilege Abuse
(ASI03) by enforcing bounded delegation depth and decaying trust across
privilege boundaries; Tripwire monitoring detects Memory and Context
Poisoning (ASI06) by alerting on unauthorized belief modifications; and
Byzantine Consensus mitigates Cascading Failures (ASI08) by requiring
supermajority agreement before collective actions, preventing a single
compromised agent from triggering system-wide degradation. This mapping
demonstrates that CIF provides a unified formal framework for threats
that OWASP currently lists as independent risks.

\textbf{NIST's Zero Trust Architecture for AI Agents} extends SP
800-207's ``never trust, always verify'\,' principles to multi-agent
environments \cite{nist2025cosais}. CIF operationalizes zero trust for
cognitive interactions: every inter-agent message is evaluated by the
firewall (continuous verification), beliefs from external sources are
sandboxed (micro-segmentation), trust scores decay exponentially with
delegation depth (least privilege), and provenance attestation provides
cryptographic message origin tracking (continuous authentication).
NIST's Control Overlays for Securing AI Systems (COSAIS) initiative,
which released its first annotated outline in January 2026 and published
a concept paper on AI agent identity and authorization in February 2026,
targets precisely the threat model that CIF formalizes---covering both
single-agent and multi-agent AI system security controls.

As these standards evolve from guidelines to compliance requirements,
CIF provides both the formal underpinning and the validated
implementation that organizations will need to demonstrate conformance.
\end{frame}

\begin{frame}{Paper Series}
\protect\phantomsection\label{paper-series}
This is Part 2 of the \emph{Cognitive Security for Multiagent Operators}
series:

\begin{itemize}
\tightlist
\item
  \textbf{Part 1: Formal Foundations} - Trust calculus, defense
  composition algebra, information-theoretic bounds
\item
  \textbf{Part 2 (This Paper): Computational Validation} -
  Implementation, attack corpus, empirical results
\item
  \textbf{Part 3: Practical Guidance} - Deployment checklists, operator
  posture, risk assessment
\end{itemize}

Together, these papers provide a complete framework for understanding,
implementing, and operating cognitive security in multiagent AI systems.
\end{frame}

\begin{frame}{Data and Code Availability}
\protect\phantomsection\label{data-and-code-availability}
The CIF implementation (defense mechanisms, evaluation framework,
analysis scripts) is available at
\url{<https://github.com/docxology/cognitive_integrity}\textgreater{}
(DOI: 10.5281/zenodo.18364128). A sanitized subset of the attack corpus
suitable for reproducibility is included; the full corpus is available
to verified researchers upon request (see \cref{sec:access-request}).
All figures, tables, and statistical analyses can be reproduced using
the provided scripts with the fixed random seed (42).
\end{frame}

\begin{frame}{Acknowledgments}
\protect\phantomsection\label{acknowledgments}
The authors thank the eight security researchers who participated in the
red team exercise, and the anonymous reviewers whose feedback
strengthened this work. We acknowledge the open-source communities
behind the multiagent frameworks evaluated in this study.
\end{frame}

\begin{frame}{Author Contributions}
\protect\phantomsection\label{author-contributions}
\textbf{Daniel Ari Friedman}: Conceptualization, Methodology, Software,
Formal analysis, Investigation, Writing -- Original Draft, Writing --
Review \& Editing, Visualization.
\end{frame}

\begin{frame}{Competing Interests}
\protect\phantomsection\label{competing-interests}
The authors declare no competing interests.
\end{frame}

\begin{frame}{Ethics Statement}
\protect\phantomsection\label{ethics-statement}
This research was reviewed and determined exempt from IRB oversight as
it did not involve human subjects. All attacks were tested against
synthetic agent configurations in sandboxed environments. Novel attack
vectors were disclosed to affected framework maintainers following a
90-day responsible disclosure policy. Dual-use risks are mitigated
through sanitization of published examples and restricted access to the
full attack corpus (see \cref{sec:dual-use}).
\end{frame}

\end{document}
